\input{LectureSlides/includes/metropolis_preamble}

% Title page info
\title[Perfect Information]{Chapter 3: Perfect-Information Games (aka Game Trees)}
\author[ECON 5001]{ECON 5001\\Prof. Paul J. Healy\\Based on: \textit{Game Theory} by Giacomo Bonanno} \color{metrop}
\institute[OSU]{}
\date[]{\vfill {\tiny Updated \today\ at\ \DTMcurrenttime}}

% ---------------------------------------------------------------------
% Trees use xgames' gametree environment. Syntax reminder:
%   \branch«overlay»[<opt>] from (<parent>) to[<edge opt>] {<children>};
%   - player=i        colors the branch and labels the parent with i's name
%   - root=w          names the parent coordinate (needed for the root)
%   - label=<angle>   places the node's player label at that angle
%   - v / h           vertical drop, horizontal spread of the children
%   - children are ";"-separated: <action>«overlay»[<opt>](<payoffs>)
%   - child i of parent (w) gets coordinate (wi), so (w22) is the second
%     child of the second child -- that's how deeper branches attach
%   - an overlay on an action, e.g. Accept<1->, highlights that edge in the
%     player's color. That is how backward-induction choices are marked
%     here, in place of Bonanno's doubled edges.
% labels=short prints "1", "2", "3" instead of full player names.
% ---------------------------------------------------------------------

\begin{document}

\frame{\maketitle}

\begin{frame}{Everything So Far Happened at Once}
In Chapter 2, players chose \textbf{simultaneously}, seeing nothing beforehand.
\bsk
But most interesting situations aren't like that:
\begin{itemize}
    \item chess
    \item an offer and a counteroffer
    \item a firm entering a market and the incumbent responding
    \item calling a play after seeing the defense line up
\end{itemize}
\vfill
Now: players move \textbf{in turn}, and everyone sees the whole history.
\vspace{-0.4in}
\end{frame}

% ==============================================================
% MOBLAB --- PLAY FIRST
% ==============================================================

% ---- MobLab run: Centipede ------------------------------------------
% Play BEFORE backward induction is defined. Default settings are fine;
% 4-6 legs, one-shot, no chat, random rematch off. Record how far the
% average pair got -- the "How Far Did We Get?" frame below wants it, and
% the "Back to the Centipede" frame near the end refers to it again.
% ---------------------------------------------------------------------
\begin{frame}{Moblab: ``Centipede Game''}
\begin{itemize}
    \item MobLab time!
    \item Two players building a firm over time
    \item As long as players invest, firm doubles in value each day
    \item Each ``turn'' one player chooses:
    \begin{itemize}
        \item \textbf{Invest}: keep the firm going.
        \item \textbf{Cash Out}: sell the firm, take 80\% of its value. Game over.
        \begin{itemize}
            \item Note: If the other person cashes out, you only get 20\%!!
        \end{itemize}
    \end{itemize}
    \item Alternating moves (P1, P2, P1, P2...)
    \item If value gets to 80, P1 must Cash Out, giving $(64,14)$
\end{itemize}
\bsk
\vfill
\centering
Feel free to talk! And play it the way you'd actually play it.
\vspace{-0.3in}
\end{frame}

\begin{frame}{How Far Did You Get?}
\vspace{0.5in}
\centering
\begin{tabular}{lc}
Most frequent \# of Turns: & \underline{\hspace{1in}} \\[0.3in]
Groups that Cashed Out in Turn 1: & \underline{\hspace{1in}} \\[0.3in]
Groups that made it to the end: & \underline{\hspace{1in}} \\[0.3in]
\end{tabular}
\end{frame}

% ==============================================================
% 3.1 TREES, FRAMES, GAMES
% ==============================================================

\begin{frame}[standout]
Trees
\end{frame}

\begin{frame}{Perfect Information}
In a \textbf{dynamic game} (or \textbf{extensive form}) players move in sequence.
\bsk
This chapter covers the case of \textbf{perfect information}:
\vfill
\centering
\large
whenever it is a player's turn to move,\\ she knows every move made so far
\vfill
\normalsize
\raggedright
\begin{itemize}
    \item \textbf{Perfect Information Games:}
    \begin{itemize}
        \item Chess, checkers, Connect Four, bargaining, sequential legislative agenda/votes
    \end{itemize}
    \item \textbf{Imperfect Information Games:}
    \begin{itemize}
        \item Stratego, Battleship, Mastermind, patent races, dynamic attacker/defender games
        \item Any game with a hidden move by ``Nature'' (see Chapter 4)
        \begin{itemize}
            \item Poker, auctions, bidding wars, diplomatic negotiations w/ private info
        \end{itemize}
    \end{itemize}
\end{itemize}
\vspace{-0.4in}
\end{frame}

\begin{frame}{Rooted Directed Trees}
We represent these games using \textbf{rooted directed trees}.
\bsk
A \textbf{rooted directed tree} is a set of \textbf{nodes} joined by directed \textbf{edges}.\\
\hspace{1in}``Directed'' edge: \textit{from} one node \textit{to} another
\bsk
Ex: Alex and Blair are out to dinner. The check arrives.
\bsk
\centering
\begin{gametree}[player names={Alex, Blair}, scale=0.82]
  \branch[player=1, root=w, label=90] from (0,0) to[v=1.4, h=5.8]
    {I've got this; Want to split?};
  \branch[player=2, label=135] from (w1) to[v=1.5, h=2.4]
    {Thanks($o_1$); Let's split($o_2$)};
  \branch[label=45] from (w2) to {Sure($o_3$); Ummm... no?($o_4$)};
\end{gametree}
\end{frame}

\begin{frame}{Rooted Directed Trees}
A \textbf{rooted directed tree} is a set of \textbf{nodes} joined by directed \textbf{edges}
\bsk
\begin{itemize}
    \item The \textbf{root} node has no edge leading into it
    \item Every other node has \textbf{exactly one} edge leading into it
    \begin{itemize}
        \item \textit{Branches out}, like a tree
        \item Contains \textit{no cycles/loops}
    \end{itemize}
    \item \textbf{Terminal nodes} have no edges going out
    \begin{itemize}
        \item They are the various game endings
        \item In a game, each has an outcome $o$ attached to it
    \end{itemize}
    \item there is a \textbf{unique path} from the root to any other node
\end{itemize}
\bsk
Set of \textbf{terminal nodes}: $Z$. All others: ``\textbf{Decision nodes}'' $D$. Then define the set of \textbf{all nodes}:
$$X = D\cup Z$$
Each node $x\in X$ represents a ``history'' of the game so far

\end{frame}

\begin{frame}{Back to the Example}
\begin{center}
\begin{gametree}[player names={Alex, Blair}, scale=0.82]
  \branch[player=1, root=w, label=90] from (0,0) to[v=1.4, h=5.8]
    {I've got this; Want to split?};
  \branch[player=2, label=135] from (w1) to[v=1.5, h=2.4]
    {Thanks($o_1$); Let's split($o_2$)};
  \branch[label=45] from (w2) to {Sure($o_3$); Ummm... no?($o_4$)};
\end{gametree}    
\end{center}
\bsk
$D=$\\
$Z=$\\
$X=$
\end{frame}

\begin{frame}{Definition: Extensive Form With Perfect Information}
A \textbf{finite extensive form with perfect information} consists of:
\bsk
\begin{itemize}
    \item A finite rooted directed tree
    \item A set of players $I=\{1,\ldots,n\}$
    \begin{itemize}
        \item Each decision node is assigned to one player
    \end{itemize}
    \item A set of actions $A$
    \begin{itemize}
        \item Each edge is assigned one action
        \item Two edges from the same node $\Rightarrow$ different actions
    \end{itemize}
    \bsk
    \item A set of outcomes $O$
    \begin{itemize}
        \item Each terminal node is assigned one outcome
    \end{itemize}
\end{itemize}
\vfill
Note: no \textbf{preferences} yet! This is a frame, not a game. Outcomes only.
\vspace{-0.3in}
\end{frame}

\begin{frame}{Dissolving the Partnership}
Nocterra Brewing Company shut down.\\
Owners: Vivian and Duncan\\
Suppose remaining assets are valued at \$100,000. How to split it?
\bsk
\begin{itemize}
    \item The charter says the senior partner, \textbf{Vivian}, proposes a split
    \item \textbf{Duncan} can Accept, or Reject and take it to court
    \item In court: each side pays \$20,000 in legal fees, and the verdict
      will 60\% to Vivian, 40\% to Duncan (no uncertainty)
    \item Vivian is allowed only two possible offers: \textbf{50-50} or \textbf{70-30}
\end{itemize}
\bsk
So rejecting yields Vivian $\$60{,}000-\$20{,}000=\$40{,}000$\\
and Duncan
$\$40{,}000-\$20{,}000=\$20{,}000$.
\vspace{-0.3in}
\end{frame}

\begin{frame}{The Frame}
\centering
\begin{gametree}[player names={Vivian, Duncan}, scale=0.85]
  \branch[player=1, root=w, label=90] from (0,0) to[v=1.3, h=5.4]
    {Offer 50-50; Offer 70-30};
  \branch[player=2, label=135] from (w1) to[v=1.5, h=2.3]
    {Accept(\makecell{$o_1$\\ \$50{,}000\\ \$50{,}000});
     Reject(\makecell{$o_2$\\ \$40{,}000\\ \$20{,}000})};
  \branch[label=45] from (w2) to
    {Accept(\makecell{$o_3$\\ \$70{,}000\\ \$30{,}000});
     Reject(\makecell{$o_4$\\ \$40{,}000\\ \$20{,}000})};
\end{gametree}
\bsk
\raggedright
{\footnotesize Top number is Vivian's (P1), bottom number is Duncan's (P2).}
\bsk
For now let's assume \textit{selfish \& greedy} preferences.
\end{frame}

\begin{frame}{Backward-Induction Reasoning}
Let's build up a new solution concept\\
\textbf{Backwards induction:} Start at the \textit{end} and work back toward the root.
\bsk
\begin{center}
    \begin{gametree}[player names={Vivian, Duncan}, scale=0.85]
  \branch[player=1, root=w, label=90] from (0,0) to[v=1.3, h=5.4]
    {Offer 50-50; Offer 70-30};
  \branch[player=2, label=135] from (w1) to[v=1.5, h=2.3]
    {Accept(\makecell{$o_1$\\ \$50{,}000\\ \$50{,}000});
     Reject(\makecell{$o_2$\\ \$40{,}000\\ \$20{,}000})};
  \branch[label=45] from (w2) to
    {Accept(\makecell{$o_3$\\ \$70{,}000\\ \$30{,}000});
     Reject(\makecell{$o_4$\\ \$40{,}000\\ \$20{,}000})};
\end{gametree}
\end{center}
\begin{itemize}
    \item<2-> Duncan accepts at both nodes
    \item<3-> Knowing this, Vivian offers 70-30
\end{itemize}
\onslide<4->{Another solution concept built on ``common knowledge of rationality''}
\end{frame}

\begin{frame}{Non-Selfish Preferences?}
What if they are \textit{not} selfish \& greedy??
\bsk
Suppose instead Duncan thinks they worked just as hard as Vivian, and that the only fair split is 50-50. They may feel strongly enough that they would give up \$10,000 to refuse an unfair offer.
\bsk
That is: \ $o_4\succ_{\text{Duncan}} o_3$.
\vfill
This will change the predicted solution:\\
Duncan Accepts 50-50 but \textit{Rejects} 70-30. So now Vivian offers 50-50.
\bsk
\centering
\large
A tree with \textit{outcomes} at the terminal nodes is a \textbf{frame}.\\
\bsk
Replace them with the players' \textit{utility values} and you have a \textbf{game}.
\vfill
\end{frame}

\begin{frame}{Adding Preferences: An Example}
Suppose Vivian is selfish and greedy: \hspace{0.35in} $o_3\succ_1 o_1\succ_1 o_2\sim_1 o_4$
\bsk
Duncan cares about fairness: \hspace{0.2in} $o_1\succ_2 o_2\sim_2 o_4\succ_2 o_3$
\bsk
\centering
\begin{tabular}{c|c|c|c|c|}
\cline{2-5}
& $o_1$ & $o_2$ & $o_3$ & $o_4$ \\
\cline{2-5}
$U_1$ (Vivian) & $2$ & $1$ & $3$ & $1$ \\
\cline{2-5}
$U_2$ (Duncan) & $3$ & $2$ & $1$ & $2$ \\
\cline{2-5}
\end{tabular}

\bsk
\raggedright
\vfill
Now replace each outcome with its payoff pair to get the actual \textbf{game}
\vspace{-0.3in}
\end{frame}

\begin{frame}{The Game}
\centering
\begin{gametree}[player names={Vivian, Duncan}, scale=0.9]
  \branch[player=1, root=w, label=90] from (0,0) to[v=1.3, h=4.6]
    {Offer 50-50<1->; Offer 70-30};
  \branch[player=2, label=135] from (w1) to[v=1.4, h=1.9]
    {Accept<1->(2,3); Reject(1,2)};
  \branch[label=45] from (w2) to {Accept(3,1); Reject<1->(1,2)};
\end{gametree}
\bsk
\raggedright
Colored edges show the backward-induction action choices.\\
Vivian now anticipates 70-30 will be rejected, offers 50-50.
\vspace{-0.2in}
\end{frame}

% ==============================================================
% 3.2 BACKWARD INDUCTION
% ==============================================================

% \begin{frame}[standout]
% Backward Induction
% \end{frame}

% \begin{frame}{The Algorithm}
% Call a node \textbf{marked} once a payoff vector is attached to it. Initially, exactly
% the terminal nodes are marked.
% \bsk
% \begin{enumerate}
%     \item Pick a decision node $x$ whose immediate successors are \textbf{all} marked.
%       Let $i$ be the player who moves at $x$.
%     \bsk
%     \item Select a choice leading to an immediate successor with the highest payoff
%       \textit{for player $i$}. Mark $x$ with that successor's whole payoff vector.
%     \bsk
%     \item Repeat until every node is marked.
% \end{enumerate}
% \vfill
% Finiteness guarantees this terminates: start at nodes followed only by terminal nodes,
% and work outward.
% \vspace{-0.3in}
% \end{frame}

\begin{frame}{Backward Induction: How to Handle Ties}
With ties, there can be multiple backwards induction predictions:
\bsk
\centering
\begin{gametree}[labels=short, scale=0.9]
  \branch[player=1, root=w, label=90] from (0,0) to[v=1.3, h=4.6] {a; b};
  \branch[player=2, label=135] from (w1) to[v=1.3, h=1.9] {c(2,1,0); d(0,0,2)};
  \branch[label=45] from (w2) to {e(3,1,0); f};
  \branch[player=3, label=45] from (w22) to[v=1.3, h=1.7] {g(1,2,1); h(0,0,1)};
\end{gametree}
\bsk
\raggedright
{\footnotesize Payoffs in order: Player 1, Player 2, Player 3.}
\vspace{-0.2in}
\end{frame}


% ==============================================================
% 3.3 STRATEGIES
% ==============================================================

\begin{frame}[standout]
Strategies
\end{frame}

\begin{frame}{Actions vs. Strategies}
A \textbf{very} important distinction: \uline{actions} vs. \uline{strategies}
\begin{center}
    \begin{gametree}[labels=short, scale=0.9]
  \branch[player=1, root=w, label=90] from (0,0) to[v=1.3, h=4.6] {a; b};
  \branch[player=2, label=135] from (w1) to[v=1.3, h=1.9] {c(2,1,0); d(0,0,2)};
  \branch[label=45] from (w2) to {e(3,1,0); f};
  \branch[player=3, label=45] from (w22) to[v=1.3, h=1.7] {g(1,2,1); h(0,0,1)};
\end{gametree}
\end{center}
\textbf{Actions:}\\
$A_1=$ \hspace{1in} $A_2=$ \hspace{2in} $A_3=$\\
\textbf{Strategies:} ``complete plans of actions''\\
$S_1=$ \hspace{1in} $S_2=$ \hspace{2in} $S_3=$
\end{frame}

\begin{frame}{A Strategy Is a Complete Plan}
\begin{block}{Definition}
A \textbf{strategy} for a player in a perfect-information game is a list of \textit{actions},
\textbf{one for each decision node of that player}.
\end{block}
\bsk
Not ``what I will do'' but ``what I \textit{would} do \uline{everywhere} I might move.''
\bsk
In the previous game, Player 2 has two decision nodes, so a strategy is a pair.
Before the game starts she doesn't know whether Player 1 will pick $a$ or $b$,
so she must plan for both:
$$(c,e)\qquad (c,f)\qquad (d,e)\qquad (d,f)$$
\vfill
Four strategies, not two.
\vspace{-0.4in}
\end{frame}

\begin{frame}{Counting Strategies}
Suppose Player 1 has three decision nodes in some game:
\begin{itemize}
    \item at the first, three choices $a_1,a_2,a_3$
    \item at the second, two choices $b_1,b_2$
    \item at the third, four choices $c_1,c_2,c_3,c_4$
\end{itemize}
A strategy fills in three blanks:
$$\Big(\underbrace{\rule{0.8in}{0.4pt}}_{\text{one of } a_1,a_2,a_3},\
\underbrace{\rule{0.8in}{0.4pt}}_{\text{one of } b_1,b_2},\
\underbrace{\rule{0.8in}{0.4pt}}_{\text{one of } c_1,\ldots,c_4}\Big)$$
\bsk
So her strategy set is a Cartesian product of the action sets she faces:
$$S_1=\{a_1,a_2,a_3\}\times\{b_1,b_2\}\times\{c_1,c_2,c_3,c_4\}$$
\centering
$3\times 2\times 4 = 24$ strategies

\raggedright
Strategy sets get large fast! Because they're a complete plan.
\vspace{-0.4in}
\end{frame}

\begin{frame}{Don't Throw Out Impossible Strategies}
Here, Player 1 moves twice:
\centering
\begin{gametree}[labels=short, scale=0.9]
  \branch[player=1, root=w, label=90] from (0,0) to[v=1.3, h=4.6] {a; b};
  \branch[player=2, label=135] from (w1) to[v=1.3, h=1.9] {c(2,1); d(0,0)};
  \branch[label=45] from (w2) to {e(3,1); f};
  \branch[player=1, label=45] from (w22) to[v=1.3, h=1.7] {g(1,2); h(1,0)};
\end{gametree}
\bsk
\raggedright
There are 2 backwards-induction predictions:
\bsk
But strategy $(a,g)$ is impossible! Playing $a$ makes $g$ irrelevant.
\bsk
But let's keep it. Player 2's choice of $f$ depended on $g$. Backwards induction requires a specification of how you play \textit{everywhere}.
\vspace{-0.2in}
\end{frame}

\begin{frame}{Other Justifications}
Other ways to justify specifying $(a,g)$ (not just $(a)$):
\bsk
\begin{enumerate}
    \item \textbf{Mistakes.} She's cautious enough to plan for the possibility that she
      intends $a$ but ends up playing $b$. \vspace{0.25in}
    \item \textbf{Delegation.} A strategy is a set of instructions handed to someone
      playing on her behalf, who needs to know what to do everywhere. \vspace{0.25in}
    \item \textbf{Beliefs.} A strategy isn't really Player 1's plan at all, but Player 2's
      \textit{belief} about what Player 1 would do. \vspace{0.25in}
\end{enumerate}
\vfill
But the main reason: \textbf{Backwards induction needs it}
\vspace{-0.3in}
\end{frame}

\begin{frame}{On Path and Off Path}
A strategy profile traces one path from the root to a terminal node.
\bsk
\begin{itemize}
    \item Nodes on that path are \textbf{on the path of play} (``\textbf{on path}'')
    \item Every other node is \textbf{off the path of play} (``\textbf{off path}'')
\end{itemize}
\bsk
A strategy must still say what to do off path. Even if certain actions aren't on-path, they make the on-path actions rational.
\vfill
\centering
\begin{gametree}[labels=short, scale=0.72]
  \branch[player=1, root=w, label=90] from (0,0) to[v=1.3, h=4.6] {\ ; \ };
  \branch[player=2, label=135] from (w1) to[v=1.3, h=1.8] {\ (\ ); \ (\ )};
  \branch[label=45] from (w2) to {\ (\ ); \ (\ )};
\end{gametree}
\vspace{-0.15in}
\end{frame}

\begin{frame}{From Tree to Table}
Every strategy profile determines a unique path, hence a unique terminal node,
hence a unique payoff vector.
\bsk
So every perfect-information game has an \textbf{associated strategic form}.
\bsk
\begin{columns}[c]
\begin{column}{0.36\textwidth}
\centering
\begin{gametree}[labels=short, scale=0.62]
  \branch[player=1, root=w, label=90] from (0,0) to[v=1.3, h=4.6] {a; b};
  \branch[player=2, label=135] from (w1) to[v=1.3, h=1.9] {c(2,1); d(0,0)};
  \branch[label=45] from (w2) to {e(3,1); f};
  \branch[player=1, label=45] from (w22) to[v=1.3, h=1.7] {g(1,2); h(1,0)};
\end{gametree}
\end{column}
\begin{column}{0.62\textwidth}
\centering
\scalebox{0.85}{%
\only<1>{\begin{matrixgame}[top={ce, cf, de, df}, left={ag, ah, bg, bh}, player names={Player 1, Player 2}]
  \payoffmatrix{2,1 & 2,1 & 0,0 & 0,0 \\ 2,1 & 2,1 & 0,0 & 0,0 \\ 3,1 & 1,2 & 3,1 & 1,2 \\ 3,1 & 1,0 & 3,1 & 1,0}
\end{matrixgame}}%
\only<2>{\begin{matrixgame}[top={ce, cf, de, df}, left={ag, ah, bg, bh}, player names={Player 1, Player 2}, br]
  \payoffmatrix{2,1 & 2,1 & 0,0 & 0,0 \\ 2,1 & 2,1 & 0,0 & 0,0 \\ 3,1 & 1,2 & 3,1 & 1,2 \\ 3,1 & 1,0 & 3,1 & 1,0}
\end{matrixgame}}}
\end{column}
\end{columns}
\bsk
\raggedright
The top two rows are identical. That's the redundancy showing up.
\vspace{-0.2in}
\end{frame}

\begin{frame}{Five Equilibria, Two Solutions}
That game has \textbf{five} Nash equilibria:
$$(ag,cf)\quad (ah,cf)\quad (bg,df)\quad (bh,ce)\quad (bh,de)$$
\bsk
It has \textbf{two} backward-induction solutions:
\begin{itemize}
    \item $\big((a,g),(c,f)\big)$ \ (because $g$ and $h$ tie for Player 1)
    \item $\big((b,h),(c,e)\big)$
\end{itemize}
\bsk
Both are Nash equilibria.\\ The other three are not backward-induction solutions.
\vfill
That gap is the whole point of the next section.
\vspace{-0.3in}
\end{frame}

\begin{frame}{Solution vs.\ Outcome}
A backward-induction \textbf{solution} is a \textit{strategy profile}. It describes
behavior everywhere, including at nodes that never get reached.
\bsk
A backward-induction \textbf{outcome} is the \textit{sequence of moves actually made}.
\bsk
\centering
\begin{tabular}{ll}
Solution $\big((a,g),(c,f)\big)$ & outcome $ac$, \ payoffs $(2,1)$ \\[0.15in]
Solution $\big((b,h),(c,e)\big)$ & outcome $be$, \ payoffs $(3,1)$ \\
\end{tabular}
\end{frame}



% ==============================================================
% 3.4 BI VS NASH
% ==============================================================

\begin{frame}[standout]
Backward Induction \& Nash
\end{frame}

\begin{frame}{The Theorem}
\textbf{Theorem.} Every backward-induction solution of a perfect-information game is a
Nash equilibrium of the associated strategic form.
\vfill
Suppose player $i$ deviates.
\begin{itemize}
    \item If she only changes what she does off the path, the play doesn't change,
      so her payoff doesn't either
    \item If she changes what she does at a node \textit{on} the path, the algorithm
      already checked that alternative and rejected it
\end{itemize}
\vfill
So backward induction is a \textbf{refinement} of Nash equilibrium: fewer solutions,
all of them Nash.
\vspace{-0.3in}
\end{frame}

\begin{frame}{Which Nash Equilibria Get Cut?}
\vfill
\centering
\large
Usually the ones resting on an \textbf{incredible threat}.
\vfill
\normalsize
\raggedright
\vspace{1.2in}
\end{frame}

\begin{frame}{The Entry Game}
IKEA has been the only big furniture store at Polaris for years, earning \$5 million.
Wayfair is deciding whether to open across I-71.
\bsk
\begin{itemize}
    \item If Wayfair stays \textbf{out}, it earns \$1 million elsewhere
    \item If Wayfair enters, IKEA can \textbf{fight}, a price war leaving both at zero
    \item Or IKEA can \textbf{accommodate} and share, and both make \$2 million
\end{itemize}
\bsk
Both firms are selfish and greedy.
\vfill
IKEA announces: ``open here and we will bury you on price.''\\
Should Wayfair believe it?
\vspace{-0.3in}
\end{frame}

\begin{frame}{The Entry Game}
\centering
\begin{gametree}[player names={Wayfair, IKEA}, scale=0.95]
  \branch[player=1, root=w, label=90] from (0,0) to[v=1.4, h=4.2]
    {Out(1,5); In<1->};
  \branch[player=2, label=0] from (w2) to[v=1.4, h=2.2]
    {Fight(0,0); Accommodate<1->(2,2)};
\end{gametree}
\bsk
\raggedright
{\footnotesize First number Wayfair, second number IKEA.}
\bsk
Backward induction: faced with entry, IKEA compares 0 to 2. It accommodates.
Knowing that, Wayfair enters.
\vspace{-0.2in}
\end{frame}

\begin{frame}{But There Are Two Nash Equilibria}
\centering
\begin{matrixgame}[top={Fight, Accommodate}, left={In, Out}, player names={Wayfair, IKEA}, br]
  \payoffmatrix{0,0 & 2,2 \\ 1,5 & 1,5}
\end{matrixgame}
\bsk
\raggedright
\begin{itemize}
    \item $(\text{In},\text{Accommodate})$, the backward-induction solution
    \item $(\text{Out},\text{Fight})$, also a Nash equilibrium!
\end{itemize}
\bsk
It survives only because Wayfair never enters, so the threat is untested.
\vfill
It is an \textbf{incredible threat}: if entry actually happened, IKEA would not
want to carry it out.
\vspace{-0.3in}
\end{frame}

\begin{frame}{Selten's Chain-Store Game}
IKEA also has a store in Cincinnati, and Wayfair could expand there next.
So the entry decision happens \textbf{one city at a time}, with everyone watching
what happened in the earlier one.
\bsk
IKEA makes the argument in CBus:
\bsk
\begin{quote}
\small
Fighting you costs me \$2 million today. But Cincy is watching. If they see me
fight, they may stay out, and then I make $\$(0+5)=\$5$ million across the two
cities. If I accommodate you here, they enter there too, and I make
$\$(2+2)=\$4$ million. So it \textit{is} in my interest to fight. Stay out.
\end{quote}
\vfill
Sounds right. Does backward induction agree?
\vspace{-0.3in}
\end{frame}

\begin{frame}{The Reply}
No, and the reason is where backward induction starts.
\bsk
\begin{itemize}
    \item Whatever happens in CBus, Cincy is the \textbf{last} interaction
    \item Nobody is watching after that, so no reputation is left to build
    \item So Cincy is just the one-shot entry game, and IKEA accommodates \vspace{0.3in}
    \item Which means Cincy goes the same way \textit{regardless} of CBus
    \item So fighting in CBus buys IKEA nothing: $\$(0+2)=\$2$ million against
      $\$(2+2)=\$4$ million
\end{itemize}
\bsk
Entry happens in both cities. IKEA accommodates both.
\vspace{-0.3in}
\end{frame}

\begin{frame}{What That Tells Us}
The reputation argument is intuitive and backward induction rejects it.
\bsk
\begin{itemize}
    \item That's not a flaw in the arithmetic. The unraveling is airtight
    \item The reputation story needs \textbf{uncertainty}: the Cincy entrant must be
      unsure what kind of opponent IKEA is
    \item In a perfect-information game, uncertainty is ruled out by definition
\end{itemize}
\vfill
We'll come back to this with better tools. For now, note the limitation.
\vspace{-0.4in}
\end{frame}

\begin{frame}{Backward Induction and IDWDS}
Backward induction resembles deleting dominated choices at each node.
\bsk
Applying IDWDS to the strategic form of a perfect-information game gives profiles
containing \textbf{at least one} backward-induction solution.
\bsk
But that's all you can say in general. The IDWDS output:
\begin{itemize}
    \item may also contain profiles that are \textbf{not} backward-induction solutions
    \item may \textbf{miss} some backward-induction solutions
\end{itemize}
\vfill
So they're related, not the same. Don't use one to compute the other.
\vspace{-0.4in}
\end{frame}

% ==============================================================
% 3.5 TWO-PLAYER GAMES
% ==============================================================

\begin{frame}[standout]
Win, Lose, or Draw
\end{frame}

\begin{frame}{Win-Lose Games}
Two players, two outcomes: Player 1 wins ($W_1$) or Player 2 wins ($W_2$).
\bsk
Player 1 strictly prefers $W_1$; Player 2 strictly prefers $W_2$. Use payoffs
$(1,0)$ and $(0,1)$.
\vfill
\centering
\textbf{Theorem.} In every finite two-player win-lose game with perfect information,
\textbf{one of the two players has a winning strategy.}
\vfill
\raggedright
A winning strategy guarantees a win \textit{no matter what the opponent does}.
\bsk
Note what it doesn't say: \textit{which} player, or what the strategy is.
\vspace{-0.3in}
\end{frame}

\begin{frame}{Let's Play}
Two players take turns choosing a number from $\{1,2,\ldots,10\}$.\\
Running total starts at 0. \textbf{First to reach 100 or more wins.}
\bsk
Player 1 moves first.
\vfill
\centering
Volunteer? I'll go second.
\vfill
\raggedright
\vspace{0.6in}
\end{frame}

\begin{frame}{Why I Was Always Going to Lose That}
The tree is far too big to draw: 10{,}000 nodes just for the first four moves.
But we can reason backward anyway.
\bsk
Call a total a \textbf{losing position} if whoever has to move from it cannot win.
\bsk
\begin{itemize}
    \item \textbf{89} is a losing position: from 89 you can reach only $90,\ldots,99$,
      and from any of those I can hit exactly 100 \vspace{0.2in}
    \item \textbf{78} is the next one down: from 78 you must reach $79,\ldots,88$,
      and from any of those I can get to 89 \vspace{0.2in}
    \item Keep going: \ $89,\ 78,\ 67,\ 56,\ 45,\ 34,\ 23,\ 12,\ \mathbf{1}$
\end{itemize}
\vspace{-0.2in}
\end{frame}

\begin{frame}{Player 1's Winning Strategy}
\vfill
\centering
\large
Start with \textbf{1}.\\
\bsk
Then every turn, choose $\mathbf{11-n}$, where $n$ is what your opponent just chose.
\vfill
\normalsize
\raggedright
Each pair of moves adds exactly 11, so the total marches $1, 12, 23, \ldots, 89, 100$.
\bsk
Notice this is a genuine \textit{strategy}: it says what to do at every position
Player 1 could face, not just the ones on the winning path.
\vspace{-0.4in}
\end{frame}

\begin{frame}{Adding a Draw}
Now three outcomes: $W_1$, $W_2$, and a draw $D$, with
$$W_1\succ_1 D\succ_1 W_2 \qquad\text{and}\qquad W_2\succ_2 D\succ_2 W_1$$
\bsk
\textbf{Theorem.} Every such game falls into exactly one of three categories:
\begin{enumerate}
    \item Player 1 has a strategy guaranteeing $W_1$
    \item Player 2 has a strategy guaranteeing $W_2$
    \item Player 1 can guarantee \textit{at least} a draw, and so can Player 2, so if
      both play well it's a draw
\end{enumerate}
\vfill
The proof is the same backward-induction argument, with three possible markings
instead of two.
\vspace{-0.3in}
\end{frame}

\begin{frame}{Which Category?}
\begin{itemize}
    \item \textbf{Tic-Tac-Toe:} category 3. Solved: perfect play draws. \vspace{0.25in}
    \item \textbf{Checkers:} category 3. Solved in 2007, after 18 years of work. \vspace{0.25in}
    \item \textbf{Chess:} \underline{\hspace{2in}} \vspace{0.25in}
\end{itemize}
\vfill
The theorem says chess is in \textit{one} of the three categories. It does not say
which, and nobody knows.
\bsk
Existence proofs are like that.
\vspace{-0.3in}
\end{frame}

% ==============================================================
% BACK TO THE CENTIPEDE
% ==============================================================

\begin{frame}{Back to the Centipede}
Apply backward induction to the game you played:
\bsk
\begin{itemize}
    \item At the \textbf{last} decision, taking beats passing. So take. \vspace{0.2in}
    \item At the second-to-last, you know the other player will take next, so passing
      gets you the smaller share. So take. \vspace{0.2in}
    \item Keep unraveling \ldots \ all the way to the first move. \vspace{0.2in}
\end{itemize}
\bsk
\centering
\large
Prediction: the very first player takes immediately.
\vfill
\normalsize
\raggedright
Our numbers said: \underline{\hspace{2.2in}}
\vspace{-0.2in}
\end{frame}

\begin{frame}{So Who's Wrong?}
Three honest possibilities:
\bsk
\begin{itemize}
    \item The \textbf{payoffs} are wrong. Players may care about the other person or
      about fairness, so the game isn't the one we wrote down \vspace{0.25in}
    \item The \textbf{reasoning} is wrong. Players may not carry the unraveling through
      that many steps \vspace{0.25in}
    \item The \textbf{beliefs} are wrong. I might reason perfectly and still pass if I
      doubt that \textit{you} will \vspace{0.25in}
\end{itemize}
\vfill
All three are live. We'll get the machinery to separate them later.
\vspace{-0.3in}
\end{frame}

% ==============================================================
% WORKED EXAMPLE
% ==============================================================

\begin{frame}[standout]
Let's Work One Out
\end{frame}

\begin{frame}{Planning Spring Break}
Ann, Ben, and Cal share a house and settle the trip in turn, each seeing every
decision made before theirs.
\bsk
\centering
\begin{gametree}[player names={Ann, Ben, Cal}, scale=0.66]
  \branch[player=1, root=w, label=90] from (0,0) to[v=1.5, h=7.4] {a; b};
  \branch[player=2, label=135] from (w1) to[v=1.5, h=3.2] {c(2,3,1); d};
  \branch[player=3, label=45] from (w2) to[v=1.5, h=3.2] {g(0,4,3); h};
  \branch[player=3, label=180] from (w12) to[v=1.5, h=1.7] {e(1,1,4); f(3,2,2)};
  \branch[player=1, label=0] from (w22) to[v=1.5, h=1.7] {i(4,0,3); j(2,2,2)};
\end{gametree}
\bsk
\raggedright
{\footnotesize Payoffs in order: Ann, Ben, Cal.}
\vspace{-0.2in}
\end{frame}

\begin{frame}{Planning Spring Break}
\begin{enumerate}
    \item How many strategies does each player have? \vspace{0.35in}
    \item Run backward induction. Where does a player face a tie? \vspace{0.35in}
    \item Write out \textbf{both} backward-induction solutions, as strategy
      profiles. \vspace{0.35in}
    \item What are the two backward-induction outcomes? \vspace{0.35in}
\end{enumerate}
\end{frame}

\begin{frame}{What Have We Learned, Charlie Brown?}
\begin{itemize}
    \item Perfect-information games are drawn as \textbf{rooted trees}; a frame becomes
      a game once preferences are added
    \item \textbf{Backward induction} works from the leaves inward; ties give several
      solutions
    \item A \textbf{strategy} covers all your decision nodes, even unreachable ones
    \item Every backward-induction solution is a \textbf{Nash equilibrium}; the extra
      Nash equilibria typically rest on \textbf{incredible threats}
    \item In finite two-player win-lose games, \textbf{somebody} has a winning strategy
\end{itemize}
\vfill
Next: what happens when players \textit{don't} see the whole history.
\vspace{-0.4in}
\end{frame}

\begin{frame}{Reading \& Practice}
\begin{itemize}
    \item Bonanno, \textit{Game Theory}, Chapter 3
    \item Exercises: \S3.7.1 (trees and frames), \S3.7.2 (backward induction),\\
      \S3.7.3 (strategies), \S3.7.4 (two-player games)
    \item Exercise 3.8 shows IDWDS missing a backward-induction solution
    \item Solutions in \S3.8
\end{itemize}
\vfill
\end{frame}

\end{document}
